% --- Sección de Desarrollo Conceptual ---
\section{Desarrollo conceptual}

\subsection{Immanuel Kant}

Immanuel Kant nacido en el oriente de prusia en 1724 en una familia de artesanos, encontró una contradicción entre la ciencia clásica, en ese entonces que era mayormente sostenido por la física mecanicista de Issac Newton en el que naturalmente determina en el tiempo y espacio una coherencia consecuente, y la moral la cual Kant sostenía que estaba fundamentalmente construida a partir de la libertad humana, es decir si las reglas del universo están regidas apartir de consecuencias entonces podríamos pensar en que las acciones humanas están estrechamente ligadas a este mundo determinista, por tanto no existiría libertad. Kant estudió meticulosamente esta brecha en su disertación, empezando estudiando la metafísica racionalista de Leibniz y la metafísica racionalista de Christian Wolff, pero los terminó descartando, porque consideró que no toda acción moral debe ser determinada por una doctrina lógica dogmática argumentando que "la moral debe ser motivada en sí misma", es decir que no podemos abstraer las máximas por argumentos sólidos lógicos, porque eso implicaría tratar de perseguir la lógica, dejando de lado la libertad.

\medskip

Al conseguir la disertación, Kant establece el Imperativo Categórico, un método por el cual se puede determinar si una máxima de un acto es moral o inmoral. Establecido por 3 principios:

\medskip

Universalidad: la máxima debe ser aplicada a todos como ley universal.

\medskip

Humanidad: la máxima debe tratar a las personas como fines y nunca considerarlos un medio para otros fines.

\medskip

Autonomía: la máxima debe de considerar a los involucrados como agentes con libre albedrío.

\medskip

Kant concluyó que la ciencia brinda leyes universales y que está sostenido por la percepción racional de lo fenoménico (Realidad consecuente), mientras que lo moral está sostenido por la libertad práctica bajo el imperativo categórico que se encuentra en la realidad nouménica.

\medskip

Para el análisis del caso es necesario considerar los 3 principios de Kant para determinar si el auto autónomo realizó un acto moral o inmoral.

\subsection{Jeremy Bentham: Utilitarismo Clásico y Cálculo Felicífico}
Bentham fundamenta el utilitarismo clásico sobre el principio de utilidad: la acción correcta es aquella que maximiza el bienestar agregado, es decir la mayor felicidad para el mayor número. La premisa es simple, la conducta humana está regida por dos estados: el dolor y el placer \citep{schofield2020jeremy}. 

Para Bentham, la moralidad es una operación aritmética denominada \textit{Cálculo Felicífico}. Este modelo propone que cualquier decisión debe evaluar variables como la intensidad, duración, certeza, proximidad  temporal, fecundidad, pureza y extensión del bienestar generado \citep{quinn2014bentham}. Lo moralmente corecta es la acción que maximiza la suma algebraica de estos factores sobre todos los individuos afectados.

Además, Bentham parte de una concepción hedonista de la naturaleza humana: placer y dolor (como se menciono más arriba), constituyen los criterios últimos de evaluación moral y legislativa. Por ello, su utilitarismo no se limita a juzgar acciones individuales, sino también leyes e instituciones, que deben ser valoradas según su capacidad para aumentar el bienestar general. En este marco, cada persona cuenta por igual dentro del cálculo moral, lo que da al utilitarismo benthamiano un carácter agregativo e imparcial. Sin embargo, esta misma pretensión de cuantificar todos los placeres en una escala común será posteriormente objeto de crítica, especialmente por parte de John Stuart Mill \citep{iepBentham}.

\subsection{John Stuart Mill: Utilitarismo Cualitativo y Justicia}
Mill evoluciona el pensamiento de Bentham al introducir una distinción cualitativa entre los placeres, argumentando que no todos los beneficios tienen el mismo peso ético \citep{west2003introduction}. Como él mismo afirma: ``Es mejor Sócrates insatisfecho que un necio satisfecho'' \citep{mill2014utilitarismo}. Esto implica que un sistema no solo debe buscar la cantidad de bienestar, sino priorizar valores superiores como la justicia y los derechos fundamentales.

Esta distinción tiene una consecuencia directa para el cálculo propuesto por Bentham: si los placeres no son cualitativamente conmensurables, no pueden sumarse en una métrica única. En consecuencia, el cálculo felicífico pierde su fundamento puramente matemático ante la complejidad de los valores cualitativos \citep{mill2014utilitarismo}.

\medskip

Mill complementa esta visión con el \textit{principio del daño} (\textit{harm principle}), principio rector que gobierna la relación entre la sociedad y el individuo. Mill establece que el único fin legítimo para interferir en la libertad de acción de una persona es la autoprotección: ``El único fin para el que el poder puede ejercerse legítimamente sobre un miembro de una comunidad civilizada, en contra de su voluntad, es evitar el daño a otros'' \citep{mill1863utilitarianism}. El bienestar propio del individuo ---físico o moral--- no es razón suficiente para coaccionarlo; sobre sí mismo, sobre su propio cuerpo y mente, el individuo es soberano.

\medskip

Este principio establece una asimetría fundamental: se puede persuadir, razonar o rogar a alguien que cambie su conducta, pero nunca compelerlo a menos que esa conducta cause daño concreto a terceros; solo el daño verificable a otros justifica la intervención \citep{mill2014utilitarismo}.

\medskip

Finalmente, Mill introduce la noción de \textit{principios secundarios}: reglas derivadas de la utilidad que actúan como árbitro cuando obligaciones morales entran en conflicto, evitando recalcular la utilidad en cada situación concreta. Como señala el propio Mill, ``si la utilidad es la fuente última de las obligaciones morales, puede invocarse para decidir entre ellas cuando sus exigencias sean incompatibles'' \citep{mill1863utilitarianism}. Esto constituye el núcleo del \textit{utilitarismo de regla}: en lugar de evaluar acto por acto, el agente ---o sistema--- sigue reglas cuya observancia generalizada maximiza el bienestar colectivo, produciendo un marco normativo estable y predecible.
